\section{Scoped obstructions inside the general frontier}
\label{sec:scoped}
The following statements refine the frontier without changing any general
exclusion count. The exact source subrings, grading supports and fixed
objects are checked separately from the general-decision certificates.

\begin{proposition}\label{prop:scoped-eight}
The rings $R8\_2\_24$, $R8\_2\_152$, $R8\_2\_153$, $R8\_2\_166$, $R8\_2\_495$, $R8\_2\_315$, $R8\_2\_317$ have no unitary
categorification. The ring $R8\_2\_312$ has no pseudo-unitary
categorification. They remain generally unresolved.
\end{proposition}
\begin{proof}
The first five rings contain the rank-four $2D2$ even ring. Its unitary
categorification is unique \cite[Theorem~3.43]{EIP}. Every invertible
bimodule over this category is an automorphism bimodule
\cite[Theorem~4.19]{Grossman}, hence is regular as a left module.
For $R8\_2\_152$, $R8\_2\_153$, $R8\_2\_166$, $R8\_2\_495$, the
odd component of a faithful $C_2$ grading has no invertible simple.
A regular module component has the regular dimension ratios, and its
equal total FP dimension fixes the scaling to one; it must contain an
invertible. This is a contradiction.

For $R8\_2\_24$, the odd component contains an invertible $u$ of
order four with $u^2=g$, where $g$ is the nontrivial neutral invertible.
Commutativity of the tensor makes $\operatorname{Ad}(u)$ fix every
neutral label, while
$\operatorname{Ad}(u)^2=\operatorname{Ad}(g)$.
The label-fixing kernel of the neutral tensor-automorphism group is
$C_2^2$: Grossman's computation gives an order-eight automorphism
group and a nontrivial $\operatorname{Ad}(g)$, and the outer,
label-fixing order-two action of \cite[Corollary~3.45]{EIP} commutes
with it. These two independent involutions exhaust the order-four
kernel. Thus no label-fixing automorphism squares to
$\operatorname{Ad}(g)$.

The rings $R8\_2\_315$, $R8\_2\_317$ contain ordinary and
pseudo-$PSU(2)_6$ subcategories with the same nontrivial invertible
$g$. In the unitary classifications their second indicators are
respectively $+1$ and $-1$. The canonical spherical structure
restricts to the canonical one on each unitary fusion subcategory,
so the indicator of the shared object cannot disagree.

Finally $R8\_2\_312$ contains a pseudo-$PSU(2)_6$ subring and an
additional $g$-fixed simple. Larson's pseudo-unitary analysis
\cite[Proposition~5.5 and proof of Theorem~5.8, Case~1]{Larson}
forces $\nu_2(g)=-1$, hence the nontrivial associator on its pointed
$C_2$ subcategory. The $g$-fixed simple gives a rank-one module over
that pointed category, forcing the associator class to be trivial by
its scalar module pentagon. This proves the pseudo-unitary
contradiction. None of the preceding arguments supplies a general
exclusion for these eight rings.
\end{proof}
